% =====================================================================
%  BENCHMARK REPORT -- glitch-token mitigation, independent evaluation
%
%  PURPOSE: source text for a larger report, in which this is one part.
%
%  HOW TO USE IT
%  -------------
%  (a) Standalone: compile as-is with pdflatex (twice, for \ref).
%  (b) As a sub-part: delete everything between BEGIN/END PREAMBLE below,
%      delete \begin{document}/\end{document} and the \maketitle block, then
%      \input{} this file. Top level is \section, so it nests under a
%      \chapter or \part without renumbering. The macros in the preamble
%      (\gp, \gc, \Ppaper, \Pcode, \Neunup, \Neundown) are used throughout
%      and must be copied across.
%
%  EVERY NUMBER is traceable. Section~\ref{sec:provenance} maps each one to a
%  file under results/ or to a page of one of the two papers. Quantities that
%  are derived rather than read (sums, set operations, re-scorings) are
%  recomputed by scripts/verify_claims.py.
%
%  Numbers in this document are for Mistral-7B-Instruct-v0.1 unless stated.
% =====================================================================

% ---------------------------------------------------- BEGIN PREAMBLE
\documentclass[11pt,a4paper]{article}
\usepackage[utf8]{inputenc}
\usepackage[T1]{fontenc}
\usepackage{lmodern}
\usepackage[margin=2.4cm]{geometry}
\usepackage{amsmath,amssymb}
\usepackage{booktabs}
\usepackage{multirow}
\usepackage{longtable}
\usepackage{graphicx}
\usepackage{xcolor}
\usepackage{caption}
\usepackage[hidelinks]{hyperref}
\usepackage{enumitem}

\newcommand{\gp}{\textsc{GlitchProber}}
\newcommand{\gc}{\textsc{GlitchCleaner}}
\newcommand{\Ppaper}{$P_{\mathrm{paper}}$}
\newcommand{\Pcode}{$P_{\mathrm{code}}$}
\newcommand{\Neunup}{\ensuremath{\mathit{Neun}_{\uparrow}}}
\newcommand{\Neundown}{\ensuremath{\mathit{Neun}_{\downarrow}}}

\title{\textbf{Independent Evaluation of Two Glitch-Token Mitigation Methods}\\[0.3em]
\large \gp{} (ASE 2024) and \gc{} (AAAI 2026)}
\author{Benchmark report}
\date{}
\begin{document}
\maketitle
{\small\tableofcontents}
% ---------------------------------------------------- END PREAMBLE

% =====================================================================
\section{Summary}
\label{sec:summary}
% =====================================================================

Two published methods claim to detect and repair \emph{glitch tokens} ---
vocabulary entries so under-trained that the model produces garbage when they
appear in a prompt. \gp{} reports an average repair rate of $50.06\%$; \gc{}
reports $86.88\%$ and positions itself as a more-than-30-point improvement. This
part of the report re-executes both under one controlled protocol on
\texttt{Mistral-7B-Instruct-v0.1}, a model both papers use, regenerating every
quantity rather than citing published values.

Five findings, ordered by how well they are supported.

\begin{enumerate}[leftmargin=*]

\item \textbf{The glitch-token population is a property of the probe, not the
model.} Holding model, decoding and vocabulary filter fixed, the same model has
$988$ glitch tokens under the prompt the papers describe and $2{,}552$ under the
prompt \gc{}'s released code actually uses. Isolating the factors, $99$--$100\%$
of that gap is the prompt \emph{wording}; the generation budget accounts for
$0.13\%$ and the correctness oracle for under $1\%$
(Section~\ref{sec:rq1}). Since a repair rate is a fraction of this population,
every published repair rate in this literature is a fraction of a
protocol-dependent denominator.

\item \textbf{That population is also a deployed component, and it breaks \gc{}'s
central design claim.} \gc{} consults the glitch set $G$ at inference time as the
gate that makes repair ``lossless'': the adapter is disabled unless the input
contains a token from $G$. But \gc{}'s own published $G$ for this model contains
\texttt{and}, \texttt{for}, \texttt{with}, \texttt{by}, \texttt{because},
\texttt{using} and the apostrophe. Measured, the gate opens on $57.4\%$ of
ordinary English sentences, on the repair evaluation's own boilerplate, and on
all $6{,}079$ examples we ran through the authors' released adapter --- it never
once closed. Losslessness is real but is not produced by the mechanism the paper
credits (Section~\ref{sec:gate}).

\item \textbf{\gc{}'s headline repair rate reproduces exactly.} Loading the
authors' own released adapter and evaluating on their own published token list
gives $2{,}408$ repaired of $2{,}539$ ($94.84\%$) against their reported $2{,}407$
($94.80\%$) --- a difference of one token (Section~\ref{sec:upstream}).

\item \textbf{\gp{}'s repair is under-determined by its own description.} Two
implementation choices the paper never states --- which tensor is adjusted, and
at which sequence positions --- move the repair rate by $7.8\times$, from
$2.23\%$ to $17.41\%$. The paper's own published figure for that configuration,
$12.92\%$, lies inside this range, so it can be neither confirmed nor refuted
(Section~\ref{sec:rq3}). Its adaptive calibration cannot be reconstructed at all:
four constants are never given.

\item \textbf{\gc{}'s reported competitor figures are rescaled, not quoted.}
\gc{} states the \gp{} results are ``taken from their original paper.'' The rates
are; the token counts are \gp{}'s \emph{rate} multiplied by \gc{}'s \emph{own}
denominator, for all five shared models (Section~\ref{sec:rescale}).

\end{enumerate}

\paragraph{On direction of error.} Most corrections this evaluation made were to
its own earlier conclusions, and most ran in the original authors' favour --- a
claimed memorisation gap, a claimed held-out shortfall, a claimed losslessness
failure, and claimed inference-speed penalties were all withdrawn after being
traced to defects in our replication rather than in the methods
(Section~\ref{sec:selfcorrection}). This is reported because it is the
characteristic failure mode of reproduction work: not missing an error, but
finding one that is your own and publishing it as someone else's.

% =====================================================================
\section{Background}
% =====================================================================

\subsection{Glitch tokens}

A language model reads \emph{tokens} from a fixed vocabulary fitted before
training. Some occur almost never in the training corpus; their embeddings are
barely updated from initialisation, and the model has no coherent representation
of what they mean. When such a token appears in a prompt the model's behaviour
degrades in ways that look like malfunction rather than ordinary error: asked to
repeat a string, it returns a different one; asked for a string's length, it
returns an unrelated number.

They matter twice over. As a reliability problem, a user who pastes text
containing one receives nonsense with no warning. As a security problem, they
occupy a region of representation space the model never learned to control, and
have been used as a lever for prompt injection and jailbreak attacks.

Both papers define a glitch token \emph{operationally}, by the same test: prompt
the model to repeat the token verbatim at temperature $0$; if it fails, the token
is glitchy. This evaluation adopts the same definition but applies it
exhaustively to the whole vocabulary, so the result is a census rather than an
estimate.

\subsection{The two methods}

\paragraph{\gp{} (Zhang et al., ASE 2024) --- detection and repair.}
Detection reads the model's internals: label a random $10\%$ of the vocabulary by
brute force, extract activation features for those tokens from a band of late
transformer layers, reduce with PCA, fit a support-vector classifier, apply it to
the remaining $90\%$, and re-validate every predicted positive with the
brute-force test. Repair edits activations at inference time: neurons reliably
active for normal tokens but under-active for a glitch token are amplified by a
factor $\beta$; neurons reliably silent for normal tokens but abnormally firing
are suppressed by a factor $\alpha$. No weights change. \textbf{No code was
released.}

\paragraph{\gc{} (Fan, Li \& Li, AAAI 2026) --- repair only.}
Attaches LoRA adapters ($<0.1\%$ additional parameters) to the MLP gate and data
projections of the same layer band and trains them on (glitch token
$\rightarrow$ correct answer) pairs. The adapters are \emph{gated}: a switch
$\lambda$ enables the branch only when the input contains a token from a known
glitch set $G$, so clean inputs are --- in the paper's account --- processed by
the mathematically unmodified base model. This gate is the basis of the paper's
``lossless'' claim. \textbf{Code, trained adapters and glitch-token lists were
released}, and are used here as the specification of record.

% =====================================================================
\section{Why these claims required verification}
\label{sec:why}
% =====================================================================

Four problems are visible in the published material \emph{before} any experiment
is run.

\subsection{The comparison spans incompatible populations}

\gc{}'s headline is a comparison against \gp{}, but the two papers worked with
different glitch-token sets for the same model. For
\texttt{Mistral-7B-Instruct-v0.1}, \gp{}'s tables imply a population of $2{,}779$
tokens; five independent derivations from its own numbers agree:
$1873/0.6741 = 2778.5$ and $1288/0.4635 = 2778.8$ and $1233/0.4437 = 2779.0$ from
detection, $359/0.1292 = 2778.6$ and $1045/0.3760 = 2779.3$ from repair. \gc{}
states $2{,}539$. A repair rate is a fraction of this population, so the two rates
are fractions of different things.

The disagreement is not specific to this model. Across the five shared models,
\gp{}'s implied population and \gc{}'s stated population differ by $7\%$ to
$35\%$, and Mistral's $9\%$ is the second smallest.

\subsection{The reported competitor counts are derived, not measured}
\label{sec:rescale}

\gc{} states that ``the \gp{} results are taken from their original paper.'' The
\emph{rates} are. The \emph{token counts} are not: for all five shared models,
\gc{}'s reported \gp{} repaired-token count equals \gp{}'s \emph{rate} multiplied
by \gc{}'s \emph{own} denominator, not the count \gp{} reports
(Table~\ref{tab:rescale}).

\begin{table}[h]
\centering
\caption{\gc{}'s Table 2 reports \gp{} token counts that are \gp{}'s published
\emph{rate} rescaled onto \gc{}'s \emph{population}. Four of five reproduce to
within $0.6$ of a token; Mistral is off by $1.3$, which the printed two-decimal
rate does not explain. The last two columns are independent corroboration.}
\label{tab:rescale}
\small
\begin{tabular}{lrrrrrr}
\toprule
Model & \gp{}'s count & \gc{} reports & \gc{}'s $|G|$ & \gp{} rate $\times$ \gc{} $|G|$ & \gp{}'s rate & \gc{}'s rate \\
\midrule
Llama-2-7b-chat          & $4{,}021$  & $2{,}968$  & $4{,}743$  & $2{,}968.2$  & $62.58\%$ & $62.58\%$ \\
Gemma-2b-it              & $13{,}638$ & $14{,}548$ & $29{,}831$ & $14{,}548.6$ & $48.77\%$ & $48.77\%$ \\
Mistral-7B-Instruct-v0.1 & $1{,}045$  & $956$      & $2{,}539$  & $954.7$      & $37.60\%$ & $\mathbf{37.65\%}$ \\
Qwen-7B-Chat             & $14{,}765$ & $13{,}320$ & $27{,}686$ & $13{,}319.7$ & $48.11\%$ & $48.11\%$ \\
Yi-6B-Chat               & $4{,}317$  & $3{,}188$  & $5{,}985$  & $3{,}187.6$  & $53.26\%$ & $\mathbf{53.27\%}$ \\
\midrule
Average                  &            &            &            &              & $50.06\%$ & $\mathbf{50.08\%}$ \\
\bottomrule
\end{tabular}
\end{table}

The two boldface disagreements are small but hard to explain any other way. Had
\gc{} quoted \gp{}'s rates, its rates would equal \gp{}'s in all five rows. They
do not, and the rows where they differ are exactly the rows where the rescaled
count required rounding: $956/2539 = 37.65\%$ and $3188/5985 = 53.27\%$. The
reported average for the \gp{} column follows the same pattern --- $50.08\%$, the
mean of the five \emph{recomputed} rates, against \gp{}'s own published $50.06\%$.
The counts were produced first and the rates derived from them.

To be fair to the authors: the rate is carried over faithfully, and rescaling is
a defensible way to place two rates on a common axis. Two things are nevertheless
wrong with it as presented. The counts are labelled as taken from the original
paper when they are computed; and the computation silently assumes \gp{}'s rate
transfers to a \emph{different set} of tokens. Section~\ref{sec:rq1} shows these
are not merely different-sized populations but different token sets, so the
assumption is substantive and untested.

\subsection{Neither paper reports variance}

\gp{} depends on a random $10\%$ vocabulary sample; \gc{} on a random adapter
initialisation and batch order. Both report single numbers. Without a spread
there is no way to distinguish a real difference from run-to-run noise --- and
Section~\ref{sec:rq2} shows the spread is large.

\subsection{\gc{} evaluates on its training data}

The paper builds a dataset from the identified glitch tokens, trains the adapters
on it, and reports repair over the glitch tokens. No held-out split is described.
If the training and evaluation populations coincide, the headline number contains
memorisation. Adding that split is the main addition this evaluation makes to
\gc{}'s own design.

% =====================================================================
\section{Method}
% =====================================================================

\subsection{Subject model and environment}

All experiments use \texttt{Mistral-7B-Instruct-v0.1}: both papers use it, it is
where the papers disagree most sharply, and it is ungated, so any reader can
reproduce this without access approvals. Single NVIDIA RTX A6000 (48\,GB), fp16,
native Windows, PyTorch 2.6 / CUDA 12.4. The original papers used $2\times$A100
(\gp{}) and an H200 (\gc{}'s speed table), so absolute wall-clock times are not
comparable and are never compared.

Restricting to one model is a deliberate depth-for-breadth trade. Every claim
below is a claim about this model; the infrastructure is model-agnostic.

\subsection{The census}

Some vocabulary entries cannot meaningfully be tested and are excluded first,
following Land \& Bartolo (2024), whom both papers adopt: control tokens; tokens
that do not decode to valid text; and \emph{unreachable} tokens, whose decoded
string re-encodes to a different id and which therefore can never arise from real
input.

The reachability test is subtler than it appears, and our first implementation
got it wrong instructively. A naive decode-then-re-encode check classified
$12{,}385$ tokens ($39\%$ of the vocabulary) as unreachable, because SentencePiece
attaches leading-space semantics to a string in isolation that it does not attach
in context. The correct procedure --- taken from \gc{}'s own released filter ---
performs the round trip behind a fixed one-token prefix, testing the token in
context. This yields $257$ filtered tokens, matching the scale both papers
report. The filtering step is load-bearing, and a plausible implementation of it
is off by two orders of magnitude.

The token under test is inserted as its \emph{raw vocabulary id}, spliced between
two pre-tokenised template halves, rather than by string concatenation followed
by re-tokenisation. This guarantees the model receives exactly the token being
tested and makes every prompt the same length, permitting exact batching.

\subsection{Two protocols, and why both are necessary}
\label{sec:protocols}

While implementing the census we found that \textbf{\gc{}'s released code does not
use the prompt its paper describes}. Two protocols therefore exist and both are
implemented.

\begin{center}
\begin{tabular}{lll}
\toprule
& \Ppaper{} & \Pcode{} \\
\midrule
Template & as written in the papers & \gc{}'s released code verbatim \\
Generation budget & $24$ tokens & $10$ tokens \\
Target string & \texttt{convert\_tokens\_to\_string} & \texttt{decode} \\
Match rule & \texttt{target.strip() in output} & \texttt{target.lstrip() in output} \\
\bottomrule
\end{tabular}
\end{center}

\Pcode{} exists so our numbers can be \emph{anchored} to the authors' published
artefacts; \Ppaper{} exists because it is what a reader of the paper would
implement. Reporting both is what separates ``the method behaves differently than
claimed'' from ``we probed it differently.'' The $24$-token budget for \Ppaper{}
is our choice --- neither paper specifies one --- and Section~\ref{sec:rq1} shows
the choice is immaterial.

\subsection{Reimplementation of \gp{}}
\label{sec:gp-reimpl}

With no public code, \gp{} was reconstructed from Algorithm 1 (detection),
Algorithm 2 (repair), Equations 7--12 and its stated hyperparameters:
$\gamma = 0.1$ sampling rate, $P = 75$ PCA components, polynomial SVM with $C=1$
and degree $3$, threshold $m=1$, key layers $19$--$28$. (The paper states that
layer band for Llama-2; we apply it to Mistral, which the paper does not
specify.)

Three aspects are \emph{under-specified in the paper} and required a choice on
our part. All three are recorded, and two are measured rather than assumed.

\begin{enumerate}[leftmargin=*]
\item \textbf{Which tensor is adjusted.} The gated MLP computes a gate signal
$\sigma(Z_1)$ and a data signal $Z_2$ and multiplies them. The paper's Figure 5
shows a ``Modified MLP gate'' \emph{and} a ``Modified MLP data'', i.e.\ the two
are adjusted separately before multiplication; we follow that reading and report
adjusting the product as an ablation.
\item \textbf{Which sequence positions are affected.} The paper never states the
position at which activations are read or the correction applied. We default to
the position of the token under test, which matches ``the activation states of
glitch tokens''; on this reading the correction fires only during prefill, since
that is the only pass in which the token is present. The alternative --- correct
at the last position of every forward pass, so the correction continues to fire
at each decoding step after the glitch token has been consumed --- is reported as
an ablation and labelled ``all positions''.
\item \textbf{The adaptive constants.} Equations 11--12 define
$\beta = k_1\Delta_\uparrow + b_1$ and $\alpha = k_2\Delta_\downarrow + b_2$, but
$k_1, b_1, k_2, b_2$ are never given; the ``adaptive process'' said to derive
them is never described; the ``range restriction'' is unspecified; and the ``set
of default values'' the text says is provided does not appear. We use identity
constants ($k_1 = k_2 = 1$, $b_1 = b_2 = 0$) with a clamp at $\alpha_{\min} = 1$.
This is \emph{our} choice and a degenerate one --- it pins $\alpha \approx 1$,
which Section~\ref{sec:rq3} independently identifies as the null-intervention
regime --- which is why the mechanism is evaluated by parameter sweep instead.
\end{enumerate}

\subsection{Reimplementation of \gc{}}

\gc{} released code, treated as the specification of record wherever the paper is
ambiguous. Our implementation matches it on LoRA rank $4$ and scaling $4$;
adapters on the MLP gate and data projections of layers $19$--$28$ only; frozen
base weights; loss on the answer span only; and the authors' exact target format
(the decoded token, unstripped, in single quotes). We adopt their optimisation
settings --- $15$ epochs, learning rate $1\times 10^{-4}$, effective batch $512$
(as $8 \times 64$ accumulation rather than their $64 \times 8$, to fit one
A6000), linear schedule with $10\%$ warmup, gradient clipping at $1.0$ --- so
that a hyperparameter mismatch cannot be offered as an explanation for any
difference in outcome.

\paragraph{The gate.} $\lambda$ is implemented as the paper and the authors' code
define it: a per-example flag, $1$ when the input contains a token from $G$ and
$0$ otherwise, multiplying the LoRA branch's contribution. The branch is still
\emph{computed} when $\lambda = 0$ --- true of the authors' implementation as
well --- so the gate removes the behavioural effect but not the arithmetic cost.

\paragraph{The held-out split (our addition).} The glitch set is split $80/20$,
the adapter trained on the $80\%$, and both halves evaluated. This separates two
things the paper's design conflates: repairing glitch tokens in general, and
memorising the specific ones shown during training.

\paragraph{Leakage control.} A token is only ``unseen'' if the adapter never
encountered its id. An id can appear inside a training \emph{sequence} without
being that sequence's subject. We scan every training sequence and exclude from
the held-out set any token whose id appears in one: $10$ of $197$ under
\Ppaper{}, $35$ of $510$ under \Pcode{}.

\subsection{Metrics, and the three denominators}
\label{sec:denominators}

Ambiguity about denominators is the specific defect this evaluation exists to
correct, so three are named here and every rate below says which it uses.

\begin{center}
\begin{tabular}{llp{7.6cm}}
\toprule
Name & Size (\Ppaper{}) & Used for \\
\midrule
census      & $988$ & the population itself; the denominator of \gp{}'s and \gc{}'s published rates \\
eval half   & $494$ & \gp{} repair. The census is split in two: one half fits the adaptive $\alpha,\beta$, both variants are scored on the other, so the comparison is not stacked in the adaptive variant's favour \\
grid sample & $500$ & each cell of the $\alpha \times \beta$ grid and the $m$-sweep, drawn from the census \\
\bottomrule
\end{tabular}
\end{center}

\noindent The eval half and the grid sample are different draws, so the same
configuration can differ slightly between tables: at $(\alpha,\beta) = (4, 1.5)$
we measure $6.48\%$ ($32/494$) and $6.80\%$ ($34/500$). That $0.3$-point
difference is the sampling noise floor for every contrast in
Section~\ref{sec:rq3}. \gc{}'s rates use its own split sizes.

\[
\text{repair rate} = \frac{\#\{\text{glitch tokens repeating correctly with the method active}\}}{|\text{stated population}|},\quad
\text{collateral} = \frac{\#\{\text{normal tokens broken}\}}{|\text{sampled normal tokens}|}
\]

Collateral breakage deserves comment: \emph{neither paper reports it}. A method
that repairs glitch tokens by degrading normal ones is not obviously an
improvement, and the quantity is cheap to measure.

% =====================================================================
\section{RQ1 --- The population is a property of the probe}
\label{sec:rq1}
% =====================================================================

\begin{table}[h]
\centering
\caption{Full-vocabulary census ($|V| = 32{,}000$). Filtering is
protocol-independent; the glitch/normal split is not. The candidate sets are
byte-identical between columns: $31{,}743$ shared ids, none exclusive to either
side.}
\label{tab:census}
\begin{tabular}{lrr}
\toprule
& \Ppaper{} & \Pcode{} \\
\midrule
Normal            & $30{,}755$ & $29{,}191$ \\
\textbf{Glitch}   & $\mathbf{988}$ & $\mathbf{2{,}552}$ \\
Unreachable       & $253$ & $253$ \\
Special           & $3$ & $3$ \\
Undecodable       & $1$ & $1$ \\
\midrule
Total             & $32{,}000$ & $32{,}000$ \\
\bottomrule
\end{tabular}
\end{table}

\paragraph{Anchoring.} Under \Pcode{} we recover $2{,}537$ of \gc{}'s $2{,}539$
published glitch tokens (Jaccard $0.9933$), with $15$ additional tokens of our
own. Our census machinery therefore reproduces the authors' published artefact
essentially exactly when run under their conditions, so the difference between
the columns of Table~\ref{tab:census} cannot be attributed to a defect in our
implementation. This anchor is what makes the rest of the section informative
rather than merely different.

\paragraph{Three populations disagree for one model.} There are now three
published or measured populations: $2{,}779$ (implied by \gp{}), $2{,}539$
(\gc{}), and ours at $988$ (\Ppaper{}) or $2{,}552$ (\Pcode{}). The first two
differ by $9\%$ while ostensibly measuring the same property of the same model.

\paragraph{Decomposing the gap.} Because the raw generations were retained, the
attribution can be settled by measurement rather than argument: greedy decoding
makes prefix truncation an exact simulation of a shorter budget, and either
oracle can be re-applied offline.

\begin{table}[h]
\centering
\caption{Attribution of the $+1{,}564$ gap, over the $31{,}743$ shared candidates.
Each factor is measured with the other two held fixed, at \emph{both} settings of
the remaining factors, because the factors interact. Token id $3$ (\texttt{NUL})
is excluded; it does not survive the CSV round trip.}
\label{tab:factorial}
\begin{tabular}{llrr}
\toprule
Factor & Measured at & $\Delta$ glitch & Share of gap \\
\midrule
Generation budget ($24 \rightarrow 10$) & either oracle & $+2$ & $0.13\%$ \\
\midrule
\multirow{2}{*}{Oracle (target derivation + match rule)}
  & \Ppaper{} generations & $+14$ & $0.90\%$ \\
  & \Pcode{} generations  & $-5$  & $-0.32\%$ \\
\midrule
\multirow{2}{*}{\textbf{Prompt wording}}
  & paper oracle  & $\mathbf{+1{,}569}$ & $\mathbf{100.3\%}$ \\
  & gccode oracle & $\mathbf{+1{,}550}$ & $\mathbf{99.1\%}$ \\
\bottomrule
\end{tabular}
\end{table}

Two things should be said plainly. First, prompt wording here is \emph{measured},
not a residual: it is the difference between two censuses that share an oracle
and a budget and differ only in the template. Second, the factors interact --- the
oracle contributes $+14$ tokens under one prompt and $-5$ under the other --- so a
single additive decomposition would be misleading, and the range is reported
instead. The conclusion is unchanged and is if anything more robust: wording
accounts for $99$--$100\%$ of the gap on either measurement path.

\paragraph{The mechanism.} Scoring both protocols with the same predicate over the
same candidates: generations consisting of newlines only number $1{,}088$ under
\Pcode{} and $0$ under \Ppaper{}; widening the predicate to whitespace-or-empty
gives $1{,}115$ against $3$. Of the $2{,}552$ tokens \Pcode{} classifies as
glitchy, $1{,}089$ ($42.7\%$) produced pure whitespace. A representative case is
token id $304$ (target \texttt{and}): under \Ppaper{} the model emits
\texttt{"' and '"} and stops; under \Pcode{} it emits ten newlines. This is
consistent with the template's trailing newline placing the model in a
newline-repetition attractor. Because three prompt edits moved together, the
effect is attributed to the wording as a whole rather than to the trailing
newline specifically; isolating it requires one further census.

\paragraph{The disagreement is bidirectional.} $146$ tokens are glitchy under
\Ppaper{} but normal under \Pcode{}, and $1{,}710$ the other way, so neither
census is a subset of the other. Any statement of the form ``$N$ tokens stop
being glitchy under protocol $X$'' must be paired with the reverse count.

\paragraph{Why this matters.} If the population is set by the probe, then so is
every published repair rate, and rates computed against different populations
cannot be compared. This is not hypothetical: \gc{}'s comparison table
(Section~\ref{sec:rescale}) reports a \gp{} token count that is \gp{}'s rate
rescaled onto \gc{}'s denominator rather than a quantity \gp{} measured.

% =====================================================================
\section{RQ2 --- \gp{} detection}
\label{sec:rq2}
% =====================================================================

\begin{table}[h]
\centering
\caption{\gp{} detection, scored against our census. Mean $\pm$ sample s.d.\ over
seeds $\{0,1,2\}$. The paper's Mistral row was measured against a population of
$2{,}779$ tokens on different hardware.}
\label{tab:detect}
\begin{tabular}{lcccc}
\toprule
& Precision & Recall & F1 & Time \\
\midrule
Ours, \Ppaper{} ($|G| = 988$)      & $0.997$ & $0.279 \pm 0.091$ & $0.431 \pm 0.116$ & $211$\,s \\
Ours, \Pcode{} ($|G| = 2{,}552$)   & $0.997$ & $0.404 \pm 0.019$ & $0.575 \pm 0.019$ & $224$\,s \\
\midrule
\gp{} claim (Mistral)              & $1.000$ & $0.674$ & $0.805$ & $42$\,m\,$39$\,s \\
\bottomrule
\end{tabular}
\end{table}

\paragraph{Recall does not reproduce, and the base rate does not account for it.}
Under \Ppaper{} recall is $0.279$. One might suspect this is depressed because our
stricter census makes positives rarer, giving the classifier less signal in its
$10\%$ training sample. Re-scoring under \Pcode{}, whose base rate ($8.0\%$)
matches the conditions the papers worked in, raises recall to $0.404$ --- so the
base rate explains part of the gap --- but it remains far from the claimed
$0.674$. The shortfall survives the control.

Because \gp{} released no code, this is reported as a failure to reproduce from
the published description, not as a refutation. Three seeds also bound variance
only loosely, and the \Ppaper{} interval alone does not exclude the paper's claim
--- it is the \Pcode{} interval that carries the inferential weight.

\paragraph{Variance is substantial and is never reported.} Under \Ppaper{} recall
across three seeds spans $0.174$--$0.335$; the only difference between those runs
is which random $10\%$ of the vocabulary trained the classifier. The interval is
much tighter under the higher base rate ($0.382$--$0.417$), which is itself
informative: the method's stability depends on how many positives its sample
happens to contain. A single reported number cannot convey this.

\paragraph{On the ``$100\%$ precision'' claim.} \gp{}'s reported $100\%$ precision
is true \emph{by construction} rather than empirically: post-validation admits a
token to the glitch set precisely when it fails the repetition test, so any
evaluation against that same oracle returns $100\%$ by definition. Scored against
an independently generated census, precision is $99.7\%$ under both protocols
($2$ and $8$ false positives across three seeds; per-seed $1,0,1$ and $3,4,1$).

The residual is a measurement artefact whose provenance matters. Our census was
generated with SDPA attention and our detection sweeps with eager attention
(required for extracting attention patterns as features); prompts, batch size,
dtype and decoding were otherwise identical. Under greedy decoding the two
kernels disagree on the generated string for $0.37$--$0.42\%$ of tokens and on
the pass/fail verdict for $0.021\%$. The divergence is systematic, not
stochastic. \emph{This attribution is an inference from the fact that two scripts
used different kernels; the controlled A/B has not been run.}

The defensible claim is narrow and, for reproducibility, consequential: a
``$100\%$ precision'' figure defined relative to a system's own oracle does not
survive contact with an independently produced ground truth, and the residual is
governed by the inference stack --- attention kernel, numerical precision,
decoding path --- none of which either paper specifies.

% =====================================================================
\section{RQ3 --- \gp{} repair}
\label{sec:rq3}
% =====================================================================

\subsection{The published calibration is unrecoverable}

\gp{}'s adaptive adjustment (Equations 9--12) cannot be reconstructed from the
paper. The functional form is given, but $k_1, b_1, k_2, b_2$ are never stated;
the ``adaptive process'' said to derive them is never described; the intervals to
which the ``range restriction'' maps $\alpha$ and $\beta$ are never given; the
``set of default values'' the text says is provided does not appear; and no
implementation was released. Algorithm 2 compounds this by computing $\beta$ and
$\alpha$ from the normal sample alone, although Equations 9--10 require
glitch-token activations, and by treating them as global scalars where
Section 4.2.2 derives them per module. The paper's own threats-to-validity
section concedes that ``the linear computation of $\alpha$ and $\beta$ proves
untenable'' while still reporting the rate obtained from it.

The consequence is precise: the reported $37.60\%$ for Mistral cannot be verified
against, or reconstructed from, the paper's description. We do not claim the
method is unimplementable --- only that the published configuration is
unrecoverable, so any reimplementation reports \emph{its own} calibration. Our
identity-constant result ($1.01\%$ under \Ppaper{}, $0.24\%$ under \Pcode{}) is
therefore \textbf{not} reported as the method's performance: that configuration
pins $\alpha \approx 1$, which the sweep below independently identifies as the
null-intervention regime, so it measures the performance of doing nothing.

\subsection{Two under-specified choices dominate the outcome}

\begin{table}[h]
\centering
\caption{Effect of the two choices the paper leaves unspecified, at the paper's
own $\alpha = 4$, $\beta = 1.5$, \Ppaper{}. All four cells are scored on the same
$494$-token evaluation half. Neuron counts are totals over the
$10 \times 2 \times 14{,}336 = 286{,}720$ candidates.}
\label{tab:implchoice}
\begin{tabular}{llccrr}
\toprule
Tensor adjusted & Positions & Repair rate & Collateral & \Neunup{} & \Neundown{} \\
\midrule
Separate gate/data (paper Fig.~5) & token only    & $6.48\%$ ($32/494$)  & $0.00\%$ & $66$  & $17{,}201$ \\
Separate gate/data                & all positions & $17.41\%$ ($86/494$) & $1.20\%$ & $638$ & $161{,}152$ \\
Product                           & token only    & $2.23\%$ ($11/494$)  & $0.00\%$ & $0$   & $45{,}353$ \\
Product                           & all positions & $15.18\%$ ($75/494$) & $1.00\%$ & $22$  & $123{,}778$ \\
\midrule
\multicolumn{2}{l}{\gp{}'s claim for this configuration} & $12.92\%$ & --- & --- & --- \\
\bottomrule
\end{tabular}
\end{table}

Three things follow.

\emph{The spread is $7.8\times$}, across four configurations that are all
defensible readings of the same paper. \emph{Where} the correction is applied
matters more than \emph{which} tensor receives it ($6.48 \rightarrow 17.41$
against $6.48 \rightarrow 2.23$), and a reader of the paper alone cannot
determine either.

\emph{The paper's own figure falls inside our range.} \gp{} reports $12.92\%$;
our four readings bracket it. The claim is therefore neither confirmed nor
refuted: it is \emph{under-determined by the paper's own description}. This is a
sharper criticism than a failure to reproduce, because it says the published
description does not pin down its own result.

\emph{Promotion is structurally starved in every reading.} The paper's key-neuron
criterion --- active above $m = 1$ in more than $99\%$ of normal tokens --- admits
between $0$ and $638$ neurons for $\beta$ to amplify, against $17{,}201$ to
$161{,}152$ for $\alpha$ to suppress, out of $286{,}720$. In the reading we take
to be the paper's, the ratio is $66 : 17{,}201$. This is a property of the
criterion the paper specifies, not of any choice we made.

\subsection{Parameter sensitivity}

\begin{table}[h]
\centering
\caption{Marginal means over the $\alpha \times \beta$ grid ($30$ cells, $500$
glitch tokens each, single seed, \Ppaper{}, separate streams, token position).}
\label{tab:grid}
\begin{tabular}{lrrrrrr}
\toprule
$\alpha$ & $1$ & $2$ & $4$ & $8$ & $16$ & spread \\
mean repair & $1.23\%$ & $4.23\%$ & $6.33\%$ & $7.37\%$ & $7.30\%$ & $6.13$\,pp \\
\midrule
$\beta$ & $0.25$ & $0.5$ & $1.0$ & $1.5$ & $2.0$ & $4.0$ \\
mean repair & $4.68\%$ & $4.64\%$ & $5.40\%$ & $5.60\%$ & $5.60\%$ & $5.84\%$ \\
\bottomrule
\multicolumn{7}{l}{\footnotesize $\beta$ marginal spread: $1.20$\,pp. Collateral
breakage is $0.0$--$0.2\%$ in every cell.}
\end{tabular}
\end{table}

The maximum anywhere in the grid is $7.60\%$ ($38/500$), attained at a
\emph{three-way tie} --- $(8, 1.0)$, $(8, 4.0)$ and $(16, 1.0)$ --- which is
itself informative: the landscape is flat enough near its top that the argmax is
not identified at this sample size. Repair rises monotonically in $\alpha$ and
saturates by $\alpha = 8$, above the paper's chosen $\alpha = 4$. The grid is
single-seed at $500$ tokens per cell, so individual cell contrasts of a few
tokens are within noise; only the marginal structure is interpreted.

\paragraph{On $\beta$.} The two factors contribute unequally --- $\alpha$'s
marginal spread is $6.13$ percentage points against $\beta$'s $1.20$ --- and
$\beta$'s effect is largest where $\alpha$ does nothing: at $\alpha = 1$, where
the suppression term is the identity, raising $\beta$ from $0.25$ to $4$ lifts
repair from $1/500$ to $13/500$. That is a real effect, and it is small: the
entire $\beta$ axis at $\alpha = 1$ moves repair by $2.4$ percentage points
against the $6.1$ points $\alpha$ moves alone. Suppression of persistently-silent
neurons carries the mechanism; promotion contributes materially only in the
regime where suppression is switched off, which is not the paper's operating
regime.

\subsection{Threshold sensitivity --- the most informative result}
\label{sec:msweep}

\begin{table}[h]
\centering
\caption{Sensitivity to the neuron-selection threshold $m$, at $\alpha = 4$,
$\beta = 1.5$, separate streams, token position, $500$ glitch and $500$ normal
tokens. Neuron counts are totals over $286{,}720$ candidates.}
\label{tab:msweep}
\begin{tabular}{lrrrr}
\toprule
$m$ & \Neunup{} & \Neundown{} & Repair rate & Collateral \\
\midrule
$0.25$        & $376$ & $530$       & $5.0\%$ ($25/500$) & $0.2\%$ \\
$0.5$         & $195$ & $2{,}222$   & $4.8\%$ ($24/500$) & $0.2\%$ \\
$1.0$ (paper) & $63$  & $17{,}041$  & $6.8\%$ ($34/500$) & $0.2\%$ \\
$2.0$         & $5$   & $150{,}755$ & $8.0\%$ ($40/500$) & $0.2\%$ \\
$4.0$         & $1$   & $270{,}882$ & $\mathbf{12.2\%}$ ($61/500$) & $0.4\%$ \\
\bottomrule
\end{tabular}
\end{table}

This isolates the mechanism by an experiment the paper does not run. Repair rises
\emph{monotonically} as the promotion set shrinks --- from $376$ neurons to a
single one --- and is highest when $\beta$ has essentially nothing to act on.
Promotion is not merely weak; removing almost all of it \emph{improves} the
outcome. What the method delivers on this model is coarse suppression of
persistently-silent neurons, and the calibrated two-factor scheme the paper
describes is not what produces the effect.

$m$ is also not a formality: moving it over a factor of $16$ changes the
selection by three orders of magnitude and the repair rate by $2.4\times$ at
nearly constant collateral cost. The paper's unjustified $m = 1$ is a
consequential unstated choice of the same kind as the tensor and the position.

\subsection{What cannot be concluded}

An earlier version of this evaluation argued that \gp{}'s claimed calibration
gain ($37.60/12.92 = 2.91\times$) was unattainable, reasoning from a repair rate
at the paper's fixed values \emph{higher} than its reported $12.92\%$. That
reasoning used a pre-correction grid in the product/all-positions configuration,
and the corrected grid reverses its direction: at the paper's own fixed values we
measure $6.80\%$, roughly \emph{half} what the paper reports. The inference
collapses with its premise and is withdrawn.

What survives is a statement about identifiability rather than attainability. The
ratio cannot be tested, for a reason the paper creates: the fixed-value baseline
it is measured against is not determined by the paper's own description --- our
four readings span $2.23$--$17.41\%$ and the published $12.92\%$ sits inside
them. A ratio between a number we cannot reconstruct and a number whose baseline
is under-determined is not a quantity a reproduction can adjudicate.

% =====================================================================
\section{RQ4 --- \gc{} repair}
\label{sec:rq4}
% =====================================================================

\subsection{Our retrained adapter}

\begin{table}[h]
\centering
\caption{\gc{} retrained with the authors' own hyperparameters, seeded, with
leaked ids excluded from the held-out set. Clean-input denominators are $500$
sampled normal tokens.}
\label{tab:gc}
\begin{tabular}{lrrccccc}
\toprule
& \multicolumn{2}{c}{split} & \multicolumn{3}{c}{glitch tokens} & \multicolumn{2}{c}{clean tokens} \\
\cmidrule(lr){2-3}\cmidrule(lr){4-6}\cmidrule(lr){7-8}
Protocol & train & held-out & train & \textbf{held-out} & adapter off & gated & forced on \\
\midrule
\Ppaper{} & $791$     & $187$ & $35.15\%$ & $\mathbf{37.43\%}$ & $0.00\%$ & $100.0\%$ & $99.6\%$ \\
\Pcode{}  & $2{,}042$ & $475$ & $82.96\%$ & $\mathbf{78.11\%}$ & $0.00\%$ & $99.6\%^{\ddagger}$ & $99.6\%$ \\
\bottomrule
\multicolumn{8}{l}{\footnotesize $^{\ddagger}$ \emph{not} a gate measurement --- see
Section~\ref{sec:gate}.}
\end{tabular}
\end{table}

The adapter-off control at $0.00\%$ confirms all repair comes from the LoRA. An
earlier version of this evaluation reported a $15$--$26$ point ``memorisation
gap'' from a run that used ad-hoc hyperparameters, an always-on adapter, unseeded
training and an unfiltered held-out set. With the authors' own hyperparameters
the gap is $4.85$ points under \Pcode{}, and under \Ppaper{} held-out
\emph{exceeds} train. Most of what had been attributed to \gc{} was our own
under-trained adapter.

That leaves an apparent shortfall --- $78.11\%$ against a published $94.80\%$ ---
and one obvious alternative explanation, namely that our training still differs
from theirs. The next subsection removes that explanation rather than arguing
about it.

\subsection{The authors' own adapter}
\label{sec:upstream}

\gc{} released its trained adapters. We load the Mistral checkpoint directly,
reproducing their \texttt{LinearWithLoRA} arithmetic
($y = \mathrm{linear}(x) + \lambda \cdot \tfrac{\alpha}{r}\,((x A) B)$, with
$A \in \mathbb{R}^{d \times r}$ and $B \in \mathbb{R}^{r \times d'}$ --- the
transpose of PEFT's convention, which is why the weights are not loaded through
PEFT) and reading rank, scaling and target layers from the checkpoint's own
configuration block rather than ours ($r = 4$, $\alpha = 4$, layers $19$--$28$;
$40$ tensors across $20$ modules). Their protocol throughout.

\begin{table}[h]
\centering
\caption{\gc{}'s released adapter, evaluated under \Pcode{}. Every row uses their
weights; the populations differ.}
\label{tab:upstream}
\begin{tabular}{lrrr}
\toprule
Population & $n$ & Repaired & Rate \\
\midrule
\textbf{Their published glitch list} (their denominator) & $2{,}539$ & $\mathbf{2{,}408}$ & $\mathbf{94.84\%}$ \\
\quad \emph{their published figure} & $2{,}539$ & \emph{$2{,}407$} & \emph{$94.80\%$} \\
\midrule
Our \Pcode{} census & $2{,}552$ & $2{,}421$ & $94.87\%$ \\
Our held-out set, restricted to their list & $473$ & $441$ & $93.23\%$ \\
Glitch tokens \emph{absent} from their list & $15$ & $15$ & $100.00\%$ \\
\midrule
Clean tokens, gate active & $500$ & $500$ & $100.00\%$ \\
Clean tokens, adapter forced on & $500$ & $500$ & $100.00\%$ \\
Control: LoRA branch zeroed & $473$ & $1$ & $0.21\%$ \\
\bottomrule
\end{tabular}
\end{table}

\paragraph{The headline claim reproduces, to within one token.} \gc{} reports
$2{,}407$ repaired of $2{,}539$; we measure $2{,}408$. The control --- the same
weights with the LoRA branch zeroed --- collapses to $0.21\%$, so the repair comes
from their parameters and not from the base model. This is a second anchor,
independent of the census anchor: our harness reproduces both \gc{}'s published
token list and its published repair rate.

\paragraph{Our shortfall was ours.} On the $473$ tokens our own adapter held out,
their adapter scores $93.23\%$ where ours scored $78.11\%$ --- same tokens, same
prompt, same decoding, same harness. The difference is a property of the adapter,
which means of our training run. The $78.11\%$-versus-$94.80\%$ comparison is
withdrawn as a statement about the method.

\paragraph{What this cannot settle.} Their adapter trained on their entire
published list, so the $473$-token row is a \emph{training-set} number for them
and says nothing about generalisation. The only genuinely unseen population is
the $15$ tokens our census finds glitchy that are absent from their list; their
adapter repairs all $15$. That is a positive signal at $n = 15$ and nothing
stronger. Whether \gc{} generalises to unseen glitch tokens therefore remains
open: our held-out design is the right instrument, and the adapter it currently
uses is not good enough for its answer to be worth quoting.

\subsection{The gate does not close}
\label{sec:gate}

\gc{}'s central design claim is that repair is \emph{lossless by construction}:
``a gating mechanism dynamically controls the activation of these branches based
on the model's input, ensuring precise intervention without disrupting normal
inference behavior.'' The gate is a set-membership test --- $\lambda = 1$ when
\emph{any} token of the input lies in $G$ --- and this is how the authors'
released code implements it
(\texttt{any(token.item() in self.glitchtokens ...)}).

The set $G$ is the same object as the census of Section~\ref{sec:rq1}, and it
inherits that section's problem in a form that is no longer presentational.
\gc{}'s own published $G$ for this model contains the vocabulary entries
\texttt{and}, \texttt{for}, \texttt{with}, \texttt{by}, \texttt{about},
\texttt{because}, \texttt{using}, \texttt{here}, the apostrophe, \texttt{=} and
\texttt{!}; $26$ of its $2{,}539$ ids fall below id $1{,}000$, i.e.\ among the
most frequent entries in the vocabulary, and $548$ decode to plain ASCII words.
Consequently:

\begin{itemize}[leftmargin=*]
\item \textbf{The gate opens on ordinary prose.} Over sentences drawn from the two
papers under study, $\lambda = 1$ for $197$ of $343$ ($57.4\%$): $62.1\%$ of
\gc{}'s own sentences and $51.4\%$ of \gp{}'s. The most frequent cause is the
token \texttt{and}, which alone opens $19.2\%$ of sentences.
\item \textbf{The gate opens on the repair evaluation's own boilerplate.} The
prompt template contains ids $304$ (\texttt{and}) and $464$ (the apostrophe),
both in $G$. So $\lambda = 1$ for every probe regardless of the token under test,
and the gate is vacuous during the headline repair measurement.
\item \textbf{It never closed in our measurements.} Under \Pcode{} our evaluation
recorded $\lambda = 1$ for all $975$ examples and $\lambda = 0$ for none,
including all $500$ clean tokens. The $99.6\%$ ``gated'' cell of
Table~\ref{tab:gc} is therefore identical to the forced-on cell by construction.
Only the \Ppaper{} row measures a closed gate, and it does so because \emph{our}
stricter census happens to exclude \texttt{and} and the apostrophe --- an accident
of our probe, not a property of the method.
\item \textbf{It never closed for the authors' adapter either.} Across all
$6{,}079$ examples of Section~\ref{sec:upstream} --- their weights, their
published $G$, their template, their membership rule --- the gate closed
\emph{zero} times.
\end{itemize}

The consequence is specific. \gc{}'s losslessness is not delivered by the
mechanism the paper credits. On ordinary input the adapter is active, so any
preservation of capability is an empirical property of a small,
narrowly-trained LoRA --- not a mathematical guarantee that the base model is what
runs. The paper's own capability table is consistent with this and, read
carefully, demonstrates it: if $\lambda$ were $0$ on benchmark prompts the \gc{}
column would be \emph{bit-identical} to the original-model column, and it is not
(Gemma-2b-it MMLU $38.14 \rightarrow 40.27$; Llama-2 GSM8K
$23.05 \rightarrow 23.88$). Those deltas are only possible if the adapter was
active.

We therefore do not dispute that \gc{} preserves capability --- their evidence
supports it and ours is consistent with it ($100\%$ of clean tokens intact with
the adapter forced on). We dispute the stated reason, and note that ``lossless by
construction'' and ``lossless in our measurements'' warrant different amounts of
confidence when the method is deployed against a different $G$.

\subsection{Coverage is bounded by a detector \gc{} does not provide}

A glitch token \emph{not} in $G$ receives $\lambda = 0$ and therefore no repair,
by construction. \gc{} is a method for repairing an enumerated list, and its
coverage in deployment is bounded by the completeness of whatever detection step
produced that list --- a step \gc{} does not perform. Composing our two
measurements: the best detection recall we could reproduce is $0.404$
(Section~\ref{sec:rq2}) and held-out repair is $0.781$ under the same protocol,
so a pipeline assembled from these two published components would repair roughly
$0.404 \times 0.781 \approx 32\%$ of the population, against a headline of
$94.80\%$. That composition assumes independence between the stages and inherits
every caveat of both, so it is an order-of-magnitude statement rather than a
measurement; it is reported because it is the number a practitioner needs and
neither paper computes it.

% =====================================================================
\section{RQ5 --- Inference cost: withdrawn}
\label{sec:rq5}
% =====================================================================

\textbf{No inference-speed result is reported.} The reason is a defect in our own
instrumentation, stated because it is the kind of error the rest of this work
exists to catch.

Two speed runs disagreed sharply --- the first gave \gp{} hooks at $0.63\times$
and the \gc{} adapter at $0.56\times$ base, the second gave both at
$\approx 0.99\times$ --- and the reversal was attributed to a configuration
change, namely that the first applied the repair correction at every decoding
step where the second applied it only at the token's own position. Re-reading the
code falsifies that account: the speed script constructs the hook object without
binding the position index, so it falls back to $-1$ and the correction fires at
every decoding step in \emph{both} runs. The file is unmodified since the
repository's first commit, so no configuration change occurred. The cause of a
$1.6\times$ swing remains unexplained, and the same script measures the \gc{}
adapter with no gate attached despite that variant having been reported as gated.

Two points nevertheless frame the eventual measurement. First, only ratios within
one environment are meaningful, since \gc{}'s published table was measured on
different hardware. Second, the gate does not exempt \gc{} from overhead: because
$\lambda$ multiplies a branch that is still computed, the adapter's arithmetic
cost is paid on clean inputs as well as glitchy ones --- true of the authors'
implementation and of ours. \gc{}'s ``negligible impact on inference speed'' is
therefore a claim about the size of the LoRA, not about the gate. We also note
that the $11.82$ tokens/s their table attributes to \gp{} is, by their own
description, their own ``simple \gp{} implementation'' rather than the original
authors' code, which was never released --- so that row measures an unspecified
reimplementation, not a published method.

% =====================================================================
\section{Threats to validity}
\label{sec:threats}
% =====================================================================

\subsection{Which findings are anchored}

\begin{itemize}[leftmargin=*]
\item \textbf{Anchored.} Everything concerning \gc{}. Our census under \Pcode{}
reproduces $2{,}537$ of their $2{,}539$ published tokens; their released code
allowed component-by-component matching of the adapter, targets and optimisation
settings; and their released adapter reproduces their headline rate to one token.
Disagreements here are informative.
\item \textbf{Not anchored.} Everything concerning \gp{}. No code was released,
and three specifications are incomplete (the adaptive constants; the adjusted
tensor; the affected positions), two of which change the outcome by an order of
magnitude. A disagreement therefore cannot be attributed with confidence, and is
reported as a failure to reproduce from the published description --- a
legitimate finding about reproducibility, not a refutation of the method.
\end{itemize}

\subsection{Construct validity: the correctness oracle is imperfect}

Both papers, and this evaluation, judge repetition by asking whether the target
string appears in the output. This is weak in two identifiable cases: tokens
whose printable form is empty, for which the test is trivially satisfied, and
very short tokens, where incidental occurrence is plausible. In our candidate set
$32$ tokens are whitespace-only and $3{,}535$ have a stripped length of at most
one character. Concretely, $3{,}253$ of the $30{,}755$ verdicts of ``normal''
under \Ppaper{} --- $10.6\%$ --- rest on a match of at most one character.

The defect is shared with the authors' released code, but sharing it is not a
defence: some fraction of every census in this literature, ours included, is
likely mislabelled for short tokens. This is the most significant unaddressed
construct-validity threat. It is orthogonal to the protocol sensitivity of
Section~\ref{sec:rq1} --- fixing the oracle would not close that gap, since the
oracle accounts for under $1\%$ of it.

\subsection{Internal validity}

One model; one adapter seed for \gc{}; a repair evaluation that consumes an
oracle glitch set rather than \gp{}'s own detector output; and an
attention-kernel attribution in Section~\ref{sec:rq2} that is an inference rather
than a controlled contrast. The oracle-set choice is deliberate --- it isolates
repair capability from detection error --- but it means the two stages are
evaluated separately rather than as the end-to-end pipeline the paper describes.

The single \gc{} seed should be stated plainly: an earlier draft attributed it to
training cost, and that was not true. One adapter trains in $149$ seconds under
\Ppaper{} and $428$ under \Pcode{}, so three seeds is a matter of minutes. It is
a gap in the work, not a constraint on it. The same applies to the three
detection seeds at $211$\,s each.

\subsection{External validity}

All results are for one 7B model. The papers' headline figures are cross-model
averages, and a single model cannot refute an average; results are stated as
claims about this model, and where a published number is compared we use that
paper's \emph{Mistral-specific} figure rather than its average.

Two considerations bear on generalisation. In \gc{}'s favour, the rescaling
finding (Section~\ref{sec:rescale}) is a five-model result requiring no compute.
Against it, the gate finding is likely to be \emph{worse} on other models, not
better: Qwen-7B-Chat's published $G$ has $44$ ids below $1{,}000$ against
Mistral's $26$, and Gemma-2b-it's $G$ covers $29{,}831$ of its vocabulary.

\subsection{Conclusion validity}

Three seeds bound variance only loosely. The $\alpha$/$\beta$ grid is single-seed
at $500$ tokens per cell, so individual cell contrasts of a few tokens are within
noise; only the marginal structure is interpreted. Where a difference is small
relative to the observed spread we say so rather than reporting it as a finding.

% =====================================================================
\section{What this evaluation corrected in itself}
\label{sec:selfcorrection}
% =====================================================================

The implementation was submitted to adversarial review before anything was
reported, and the results were then re-checked against the stored artefacts a
second time by a different author. Both passes found errors. They are reported
because the difference between the pre- and post-correction numbers measures how
sensitive these methods are to choices the papers do not pin down --- and because
a benchmark's authority rests on its auditability.

\begin{table}[h]
\centering
\caption{Corrections, and their direction. ``Favours'' records whether the
corrected version is better or worse for the original authors.}
\label{tab:corrections}
\small
\begin{tabular}{p{5.2cm}p{4.2cm}p{4.2cm}l}
\toprule
Claim as first written & What was wrong & Corrected form & Favours \\
\midrule
$15$--$26$ point memorisation gap in \gc{} & our training: 3 epochs not 15, wrong LR, unseeded, dropped gradients & gap is $4.85$ points; held-out exceeds train under \Ppaper{} & authors \\
\addlinespace
Held-out $78.11\%$ vs claimed $94.80\%$ is a reproduction shortfall & still our training: their adapter scores $93.23\%$ on the same tokens & the claim reproduces to one token & authors \\
\addlinespace
Losslessness is supported by the gate & the gate never closed; $99.6\%$ was the always-on number relabelled & losslessness is real but not from the gate & mixed \\
\addlinespace
$37$--$44\%$ inference-speed penalties & v1 configuration artefact & withdrawn & authors \\
\addlinespace
$\approx 1\%$ inference-speed overhead & the script never binds the position index; both runs measured the same configuration & withdrawn; no result reported & neither \\
\addlinespace
\gp{}'s calibration gain is unattainable & argued from the configuration we classify as our own error; corrected value is below, not above, the paper's & withdrawn; replaced by an identifiability argument & authors \\
\addlinespace
``$\beta$ is inert'' & partly an artefact of adjusting the wrong tensor & $\beta$ is structurally starved, not inert; $m$-sweep establishes the point independently & mixed \\
\addlinespace
Census gap is due to generation length & length accounts for $0.13\%$; wording for $99$--$100\%$ & umbrella claim survives and strengthens; stated cause reversed & neither \\
\addlinespace
Factors do not interact & the oracle effect changes sign between prompts & report the range & neither \\
\addlinespace
$100\%$ precision fails due to fp16 non-determinism & the divergence is a deterministic kernel mismatch between our own scripts & narrower claim about the inference stack & authors \\
\bottomrule
\end{tabular}
\end{table}

Two further errors of \emph{implementation} were found and fixed before any of
the above: \gp{} repair adjusted the already-multiplied product of the gate and
data streams rather than the two streams separately; and \gc{}'s $\lambda$ gate
was absent entirely, so every claim about clean-input behaviour was measured on
an always-on adapter. Nine smaller correctness defects were also fixed (unseeded
training; a dropped final gradient-accumulation window each epoch; adaptive
parameters fitted on their own evaluation set; an absolute value inserted into a
published signed ratio; held-out leakage; a training-target format mismatch; an
inconsistent attention kernel; a stale checkpoint silently resumed; and an
advertised experiment that had never been executed).

\paragraph{The lesson is not ``audit the implementation''.} The first pass found
errors in the \emph{code}. The two largest later corrections --- losslessness and
inference speed --- are errors in the \emph{reporting}: the code did something
defensible and the prose described something else, and every number involved was
real, stored and reproducible. Auditing an implementation does not catch these.
What caught them was re-reading each stored artefact against the sentence written
about it.

\paragraph{The safeguards that worked were structural, not diligence.} Retained
raw generations made the factorial decomposition possible after the fact;
recorded provenance --- including a dirty-working-tree flag --- exposed artefacts
attributed to commits that could not have produced them; a configuration
fingerprint caught a ``corrected'' parameter grid that had silently resumed its
own pre-correction checkpoint; and a reviewer instructed to disbelieve the
documentation found what internal review had not.

% =====================================================================
\section{Recommendations}
% =====================================================================

\begin{enumerate}[leftmargin=*]
\item \textbf{Publish the census and the probe, not just the rate.} The population
a repair rate is a fraction of is set by the prompt used to find it, and $99\%$ of
a $2.5\times$ difference in that population comes from wording alone. A paper
reporting a repair rate without its census and its verbatim probe is not
reporting a reproducible quantity. The same applies to reporting a competitor's
rate: rescaling it onto one's own denominator produces a number the competitor
never measured.
\item \textbf{Do not let the census become a runtime component without saying
so.} \gc{}'s gate consults the glitch set at inference time, which converts a
measurement artefact into a behavioural one. If a set built by a probe is going
to gate a production code path, its false-positive rate on ordinary text is a
first-order property and should be reported.
\item \textbf{Evaluate learned repair on held-out tokens.} A method trained on an
enumerated list and evaluated on that same list reports a quantity that includes
memorisation. The split costs nothing and changes the interpretation.
\item \textbf{Report collateral damage.} Neither paper states what its repair does
to normal tokens. It is cheap to measure and is the natural counterweight to a
repair rate.
\item \textbf{Specify the intervention completely.} For \gp{}, the withheld
constants are the visible gap, but the unstated tensor and unstated positions
change the outcome by an order of magnitude and are invisible to a reader. A
method description should be sufficient to determine what the code does.
\item \textbf{Specify the inference stack.} Attention kernel, precision and
decoding path change the oracle's verdict on a measurable fraction of tokens.
Where a claim is $100\%$ by construction, that should be stated as a definition
rather than presented as a measurement.
\end{enumerate}

% =====================================================================
\section{Conclusion}
% =====================================================================

\paragraph{What to believe.} \gc{}'s headline repair rate is real: the authors'
own released adapter reproduces it to within one token. Two attempts to find a
shortfall in it --- a memorisation argument and a held-out comparison --- both
turned out to be shortfalls in our replication of their training. It remains
measured on the training population, so it is not a generalisation estimate, and
the one unseen population we could construct ($n = 15$) is too small to settle
that. \gp{}'s repair mechanism works but weakly, and its published configuration
cannot be reconstructed at all.

\paragraph{What to deploy.} Three things bound what either method delivers, none
visible from the papers. Coverage is bounded by detection: \gc{} repairs only
tokens on a list it is given, and the best detection recall we could reproduce is
$0.404$, so a pipeline built from these two published components covers roughly a
third of the population, not $94.80\%$. The ``lossless'' guarantee is not a
guarantee: the gate that provides it opens on the majority of ordinary text.
And the population itself is a property of the probe --- anyone deploying these
methods will build their own glitch set with their own prompt, and a $2.5\times$
difference in that set is attainable by wording alone.

\paragraph{What remains open.} \gp{}'s adaptive calibration is unrecoverable, so
its principal claim can be neither reproduced nor refuted. Its fixed-value
baseline is under-determined by its own description. Whether \gc{} generalises to
unseen glitch tokens is unresolved: the design that answers it is ours, but our
adapter is not yet good enough for its answer to be worth quoting, and the
authors' adapter has seen every token on their list. And the inference-speed
measurement was withdrawn rather than published.

\paragraph{On direction.} Of the substantive corrections made here, most moved
\emph{in the original authors' favour}. What survived in the other direction is
narrower and more specific: the census is a property of the probe, the competitor
counts are rescaled rather than quoted, \gp{}'s repair is under-specified to the
point of being unfalsifiable, and \gc{}'s losslessness is credited to the wrong
mechanism. That asymmetry is the most useful thing to report about how
reproduction studies go wrong: the easy failure mode is not missing an error, it
is finding one that is your own and publishing it as someone else's.

% =====================================================================
\section{Provenance}
\label{sec:provenance}
% =====================================================================

Every number above is either produced by code in this repository and stored under
\texttt{results/}, or read from one of the two papers. The two are distinguished
below, because conflating them is the failure this work is about. Derived
quantities --- sums, set operations, re-scorings, none of which is the direct
output of any run --- are recomputed by \texttt{scripts/verify\_claims.py}, which
prints every figure in Sections~\ref{sec:rq1}, \ref{sec:rq3} and \ref{sec:gate}.

\begin{longtable}{p{4.4cm}p{4.6cm}p{5.0cm}}
\toprule
\textbf{Result} & \textbf{Script} & \textbf{Artefact} \\
\midrule
\endhead
Census (Table~\ref{tab:census}) & \texttt{run\_ground\_truth.py} & \texttt{results/ground\_truth/<m>/\{tokens.csv, summary.json, sweep\_checkpoint.csv\}} \\
\Pcode{} census & \texttt{-{}-protocol gccode} & \texttt{.../gccode/} \\
Budget cell (24 tokens) & \texttt{-{}-max-new-tokens 24 -{}-tag budget24} & \texttt{.../gccode/budget24/} \\
Factorial (Table~\ref{tab:factorial}), anchoring, mechanism, oracle counts & \texttt{verify\_claims.py} & derived from the \texttt{sweep\_checkpoint.csv} files \\
Detection (Table~\ref{tab:detect}) & \texttt{run\_gp\_detect.py} & \texttt{results/gp\_detect/<m>/\{runs.csv, summary.json\}} \\
Implementation 2$\times$2 (Table~\ref{tab:implchoice}) & \texttt{run\_gp\_repair.py} with \texttt{-{}-stream-mode}, \texttt{-{}-position} & \texttt{results/gp\_repair/<m>/\{summary.json, ablation\_product/, abl\_separate\_last/, abl\_product\_last/\}} \\
Grid (Table~\ref{tab:grid}), $m$-sweep (Table~\ref{tab:msweep}) & \texttt{run\_gp\_alpha\_beta\_sweep.py} & \texttt{.../\{alpha\_beta\_grid.csv, m\_sweep.csv, m\_sweep\_meta.json\}} \\
Pre-correction grid & (retained for contrast) & \texttt{.../alpha\_beta\_grid.v1\_product\_last.csv} \\
Our adapter (Table~\ref{tab:gc}) & \texttt{run\_gc\_train.py}, \texttt{run\_gc\_eval.py} & \texttt{results/gc/<m>/\{split.json, train\_meta.json, eval.json\}} \\
Their adapter (Table~\ref{tab:upstream}) & \texttt{run\_gc\_upstream\_eval.py} & \texttt{results/gc\_upstream/<m>/gccode/eval.json} \\
Gate statistics (Section~\ref{sec:gate}) & \texttt{verify\_claims.py} & \texttt{eval.json:gate\_stats} + the tokenizer + \gc{}'s published list \\
Inference speed & \texttt{run\_speed.py} & \emph{withdrawn --- see Section~\ref{sec:rq5}} \\
Run narrative, with corrections & --- & \texttt{docs/RUNLOG.md} \\
\midrule
\multicolumn{3}{l}{\emph{Read from the papers, not measured here:}} \\
Table~\ref{tab:rescale}; implied populations & --- & \gc{} Table 2; \gp{} Tables 3 and 5 \\
\gp{}'s claimed detection row & --- & \gp{} Table 3 \\
\gc{}'s capability deltas (Section~\ref{sec:gate}) & --- & \gc{} Table 4 \\
\gc{}'s published glitch sets and adapters & --- & \texttt{third\_party/GlitchCleaner/} \\
\bottomrule
\end{longtable}

\paragraph{One figure is not our own measurement.} The statement in
Section~\ref{sec:gate} that the gate opens on ordinary benchmark text is
supported here by our own sentence-level measurement ($57.4\%$ over the two
papers' text). A prior adversarial review of this project additionally reported
$92.27\%$ of GSM8K items and $89.77\%$ of MMLU items opening the gate; those
figures are recorded in \texttt{docs/evidence\_adjudications/upstream\_code.md}
and have \emph{not} been independently re-measured here. They are cited as
corroboration, not as a result of this evaluation.

\end{document}
